%%%%%%%%%%%%
%
% $Autor: Wings $
% $Datum: 2019-03-05 08:03:15Z $
% $Pfad: TemplateSensor $
% $Version: 4250 $
% !TeX spellcheck = en_GB/de_DE
% !TeX encoding = utf8
% !TeX root = filename 
% !TeX TXS-program:bibliography = txs:///biber
%
%%%%%%%%%%%%

% Structure
chapter{Mikrocontroller}

\chapter{General}


\chapter{Mikrokontroller}


\subsection{Aufbau Mikrokontroller}	


\subsection{Funktionsprinzip}


\chapter{Arduino Uno R4 WiFi}


\subsection{Technische Daten}
Nachfolgende Tabelle (	) zeigt die im Datenblatt \cite{Arduino:2026} entnommen technischen Eigenschaften des Arduino's.

\begin{table}[h!]
	\caption{Verbauliste Arduino Uno R4 WiFi \cite{Arduino:2026}}
	\begin{center}
		\begin{tabular}{|l|l|l|}
			\hline
			\textbf{Parameter} &\textbf{Detail} & \textbf{Referenzierung}\\
			\hline
			Mikrocontroller & R7FA4M1AB3CFM0 AA0 & U1	\\
			\hline
			Microplexer & NLASB3157DFT2G & U2	\\	
			\hline
			Abwärtswandler & ISL854102FRZ-T & U3	\\	
			\hline
			Logikpegelwandler (5 V – 3,3 V) & TXB0108DQSR  & U4	\\	
			\hline
			3,3-V-Linearregler & SGM2205-3.3XKC3G/TR & U5	\\	
			\hline
			Multiplexer & NLASB3157DFT2G & U6	\\	
			\hline
			LED-Matrix  & 12x8 Rot & U-LEDMATRIX\\	
			\hline
			Mikrocontroller & ESP32-S3-MINI-1-N8 & M1\\	
			\hline
			RESET-Taste & ------- & PB1\\	
			\hline
			Analog-Steckleiste & Ein-Ausgaenge & JANALOG\\	
			\hline
			Digital-Steckleiste & Ein-Ausgaenge & JANADIGITAL\\
			\hline
			VRTC-Steckleiste & --------- & JOFF	\\	
			\hline
			USB-C-Anschluss & CX90B-16P & J1	\\	
			\hline
			Qwiic-Anschluss & I2C-Kommunikationsprotokoll & J2	\\	
			\hline
			ICSP-Steckleiste  & --------- & J3	\\	
			\hline
			DC-Buchse & --------- & J5	\\	
			\hline
			ESP-Steckleiste & --------- & J6	\\	
			\hline
			LED TX & serielle Übertragung & DL1	\\	
			\hline
			LED RX & serieller Empfang & DL2	\\	
			\hline
			LED Power & gruen & DL3	\\	
			\hline
			LED SCK & Serieller Takt & DL4	\\	
			\hline
			Schottky-Diode & PMEG6020AELRX  & D1	\\	
			\hline
			Schottky-Diode & PMEG6020AELRX  & D2	\\
			\hline
			ESD-Schutz & PRTR5V0U2X,215 & D3	\\		
			\hline
		\end{tabular}
		\label{XXXXXXXXXXXXXXXXXXXXX}
	\end{center}
\end{table} 
Nachfolgend wird auf den Microcontroller

\subsection{Handbuch}


\subsection{Inbetriebnahme}


\subsection{Anschlussbeispiel}



\subsection{Code}


\chapter{Anwendungsbeschreibung}



\chapter{Tests}


\chapter{Zusatzinformationen}